\begin{name}
	{\tenchude}
	{TOÁN 10 - BPT, HỆ BPT BẬC NHẤT HAI ẨN}
	{LỚP TOÁN THẦY PHÁT}
	{---...---}
\end{name}

\Opensolutionfile{ans}[ans/ans-0D2-B1-MienNghiem]

%%%%%%===================================================================%%%%%%
%%%%%%============================= DẠNG 1 ===============================%%%%%%
%%%%%% Bờ là đường thẳng song song với trục tung: x = m
%%%%%%===================================================================%%%%%%
% \subsubsection{Bờ là đường thẳng $x=m$}

%%%==============Câu_1 — x = 0 (trục tung), tô bên trái, bờ NÉT LIỀN=============%%%
\begin{ex}%[VNPmath - tự biên soạn]%[GV. Vũ Ngọc Phát]%[0D2N1-2]
	\immini[thm]{Phần không bị gạch trong hình vẽ là miền nghiệm của bất phương trình nào sau đây?
		\choice
		{$x\leq 0$}
		{\True $x\geq 0$}
		{$y\leq 0$}
		{$y\geq 0$}}
	{\begin{tikzpicture}[scale=.7,thick,>=stealth,declare function={xmin=-3;xmax=3;ymin=-2;ymax=3;m=0;}]
			\fill[pattern=north east lines] (xmin,ymin) rectangle (xmax,ymax);
			\fill[white] (m,ymin) rectangle (xmax,ymax);
			\draw[->] (xmin,0) -- (xmax,0)node[above]{$x$};
			\foreach \x in {-2,-1,1,2,3}
			\draw[shift={(\x,0)}] (0pt,2pt) -- (0pt,-2pt) node[below] {\footnotesize $\x$};
			\draw[->,color=black] (0,ymin) -- (0,ymax)node[left]{$y$};
			\foreach \y in {-1,1,2}
			\draw[shift={(0,\y)}] (2pt,0pt) -- (-2pt,0pt) node[left] {\footnotesize $\y$};
			\node[above left] at (0,0){$O$};
		\end{tikzpicture}}
	\loigiai{
		Đường thẳng bờ là trục tung $d\colon x=0$ (vẽ nét liền nên kể cả bờ).\\
		Lấy điểm $M(1;0)$ thuộc phần không bị gạch, ta có $1\geq 0$ nên phần không bị gạch là miền nghiệm của bất phương trình $x\geq 0$.
	}
\end{ex}

%%%==============Câu_2 — x = 1, tô bên trái, bờ NÉT LIỀN=============%%%
\begin{ex}%[VNPmath - tự biên soạn]%[GV. Vũ Ngọc Phát]%[0D2N1-2]
	\immini[thm]{Phần không bị gạch trong hình vẽ là miền nghiệm của bất phương trình nào sau đây?
		\choice
		{\True $x\geq 1$}
		{$x\leq 1$}
		{$x\geq -1$}
		{$x\leq -1$}}
	{\begin{tikzpicture}[scale=.7,thick,>=stealth,declare function={xmin=-2;xmax=4;ymin=-2;ymax=3;m=1;}]
			\fill[pattern=north east lines] (xmin,ymin) rectangle (xmax,ymax);
			\fill[white] (m,ymin) rectangle (xmax,ymax);
			\draw[thick] (m,ymin) -- (m,ymax);
			\draw[->] (xmin,0) -- (xmax,0)node[above]{$x$};
			\foreach \x in {-1,1,2,3,4}
			\draw[shift={(\x,0)}] (0pt,2pt) -- (0pt,-2pt) node[below] {\footnotesize $\x$};
			\draw[->,color=black] (0,ymin) -- (0,ymax)node[left]{$y$};
			\foreach \y in {-1,1,2}
			\draw[shift={(0,\y)}] (2pt,0pt) -- (-2pt,0pt) node[left] {\footnotesize $\y$};
			\node[above left] at (0,0){$O$};
		\end{tikzpicture}}
	\loigiai{
		Đường thẳng bờ $d\colon x=1$ (vẽ nét liền nên kể cả bờ).\\
		Lấy điểm $M(2;0)$ thuộc phần không bị gạch, ta có $2\geq 1$ nên phần không bị gạch là miền nghiệm của bất phương trình $x\geq 1$.
	}
\end{ex}

%%%==============Câu_3 — x = -1, tô bên phải, bờ NÉT LIỀN=============%%%
\begin{ex}%[VNPmath - tự biên soạn]%[GV. Vũ Ngọc Phát]%[0D2N1-2]
	\immini[thm]{Phần không bị gạch trong hình vẽ là miền nghiệm của bất phương trình nào sau đây?
		\choice
		{$x\geq -1$}
		{$x\geq 1$}
		{\True $x\leq -1$}
		{$x\leq 1$}}
	{\begin{tikzpicture}[scale=.7,thick,>=stealth,declare function={xmin=-4;xmax=2;ymin=-2;ymax=3;m=-1;}]
			\fill[pattern=north east lines] (xmin,ymin) rectangle (xmax,ymax);
			\fill[white] (xmin,ymin) rectangle (m,ymax);
			\draw[thick] (m,ymin) -- (m,ymax);
			\draw[->] (xmin,0) -- (xmax,0)node[above]{$x$};
			\foreach \x in {-3,-2,-1,1,2}
			\draw[shift={(\x,0)}] (0pt,2pt) -- (0pt,-2pt) node[below] {\footnotesize $\x$};
			\draw[->,color=black] (0,ymin) -- (0,ymax)node[left]{$y$};
			\foreach \y in {-1,1,2}
			\draw[shift={(0,\y)}] (2pt,0pt) -- (-2pt,0pt) node[left] {\footnotesize $\y$};
			\node[above left] at (0,0){$O$};
		\end{tikzpicture}}
	\loigiai{
		Đường thẳng bờ $d\colon x=-1$ (vẽ nét liền nên kể cả bờ).\\
		Lấy điểm $M(-2;0)$ thuộc phần không bị gạch, ta có $-2\leq -1$ nên phần không bị gạch là miền nghiệm của bất phương trình $x\leq -1$.
	}
\end{ex}

%%%==============Câu_4 — x = 2, tô bên phải, bờ NÉT ĐỨT=============%%%
\begin{ex}%[VNPmath - tự biên soạn]%[GV. Vũ Ngọc Phát]%[0D2N1-2]
	\immini[thm]{Phần không bị gạch (không kể đường thẳng $d$) trong hình vẽ là miền nghiệm của bất phương trình nào sau đây?
		\choice
		{$x>2$}
		{$x<3$}
		{$x>3$}
		{\True $x<2$}}
	{\begin{tikzpicture}[scale=.7,thick,>=stealth,declare function={xmin=-1;xmax=5;ymin=-2;ymax=3;m=2;}]
			\fill[pattern=north east lines] (xmin,ymin) rectangle (xmax,ymax);
			\fill[white] (xmin,ymin) rectangle (m,ymax);
			\draw[thick,dashed] (m,ymin) -- (m,ymax);
			\draw[->] (xmin,0) -- (xmax,0)node[above]{$x$};
			\foreach \x in {-1,1,2,3,4,5}
			\draw[shift={(\x,0)}] (0pt,2pt) -- (0pt,-2pt) node[below] {\footnotesize $\x$};
			\draw[->,color=black] (0,ymin) -- (0,ymax)node[left]{$y$};
			\foreach \y in {-1,1,2}
			\draw[shift={(0,\y)}] (2pt,0pt) -- (-2pt,0pt) node[left] {\footnotesize $\y$};
			\node[above left] at (0,0){$O$};
		\end{tikzpicture}}
	\loigiai{
		Đường thẳng bờ $d\colon x=2$ vẽ nét đứt nên không kể bờ.\\
		Lấy điểm $O(0;0)$ thuộc phần không bị gạch, ta có $0<2$ nên phần không bị gạch là miền nghiệm của bất phương trình $x<2$.
	}
\end{ex}

%%%==============Câu_5 — x = -2, tô bên trái, bờ NÉT ĐỨT=============%%%
\begin{ex}%[VNPmath - tự biên soạn]%[GV. Vũ Ngọc Phát]%[0D2N1-2]
	\immini[thm]{Phần không bị gạch (không kể đường thẳng $d$) trong hình vẽ là miền nghiệm của bất phương trình nào sau đây?
		\choice
		{$x<-3$}
		{$x>-3$}
		{$x<-2$}
		{\True $x>-2$}}
	{\begin{tikzpicture}[scale=.7,thick,>=stealth,declare function={xmin=-5;xmax=1;ymin=-2;ymax=3;m=-2;}]
			\fill[pattern=north east lines] (xmin,ymin) rectangle (xmax,ymax);
			\fill[white] (m,ymin) rectangle (xmax,ymax);
			\draw[thick,dashed] (m,ymin) -- (m,ymax);
			\draw[->] (xmin,0) -- (xmax,0)node[above]{$x$};
			\foreach \x in {-4,-3,-2,-1,1}
			\draw[shift={(\x,0)}] (0pt,2pt) -- (0pt,-2pt) node[below] {\footnotesize $\x$};
			\draw[->,color=black] (0,ymin) -- (0,ymax)node[left]{$y$};
			\foreach \y in {-1,1,2}
			\draw[shift={(0,\y)}] (2pt,0pt) -- (-2pt,0pt) node[left] {\footnotesize $\y$};
			\node[above left] at (0,0){$O$};
		\end{tikzpicture}}
	\loigiai{
		Đường thẳng bờ $d\colon x=-2$ vẽ nét đứt nên không kể bờ.\\
		Lấy điểm $O(0;0)$ thuộc phần không bị gạch, ta có $0>-2$ nên phần không bị gạch là miền nghiệm của bất phương trình $x>-2$.
	}
\end{ex}

%%%%%%===================================================================%%%%%%
%%%%%%============================= DẠNG 2 ===============================%%%%%%
%%%%%% Bờ là đường thẳng song song với trục hoành: y = m
%%%%%%===================================================================%%%%%%
% \subsubsection{Bờ là đường thẳng $y=m$}

%%%==============Câu_6 — y = 0 (trục hoành), tô phía dưới, bờ NÉT LIỀN=============%%%
\begin{ex}%[VNPmath - tự biên soạn]%[GV. Vũ Ngọc Phát]%[0D2N1-2]
	\immini[thm]{Phần không bị gạch trong hình vẽ là miền nghiệm của bất phương trình nào sau đây?
		\choice
		{$y\leq 0$}
		{$x\leq 0$}
		{$x\geq 0$}
		{\True $y\geq 0$}}
	{\begin{tikzpicture}[scale=.7,thick,>=stealth,declare function={xmin=-2;xmax=4;ymin=-2.5;ymax=2.5;m=0;}]
			\fill[pattern=north east lines] (xmin,ymin) rectangle (xmax,ymax);
			\fill[white] (xmin,m) rectangle (xmax,ymax);
			\draw[->] (xmin,0) -- (xmax,0)node[above]{$x$};
			\foreach \x in {-1,1,2,3,4}
			\draw[shift={(\x,0)}] (0pt,2pt) -- (0pt,-2pt) node[below] {\footnotesize $\x$};
			\draw[->,color=black] (0,ymin) -- (0,ymax)node[left]{$y$};
			\foreach \y in {-2,-1,1,2}
			\draw[shift={(0,\y)}] (2pt,0pt) -- (-2pt,0pt) node[left] {\footnotesize $\y$};
			\node[above left] at (0,0){$O$};
		\end{tikzpicture}}
	\loigiai{
		Đường thẳng bờ là trục hoành $d\colon y=0$ (vẽ nét liền nên kể cả bờ).\\
		Lấy điểm $M(0;1)$ thuộc phần không bị gạch, ta có $1\geq 0$ nên phần không bị gạch là miền nghiệm của bất phương trình $y\geq 0$.
	}
\end{ex}

%%%==============Câu_7 — y = 1, tô phía dưới, bờ NÉT ĐỨT=============%%%
\begin{ex}%[VNPmath - tự biên soạn]%[GV. Vũ Ngọc Phát]%[0D2N1-2]
	\immini[thm]{Phần không bị gạch (không kể đường thẳng $d$) trong hình vẽ là miền nghiệm của bất phương trình nào sau đây?
		\choice
		{\True $y>1$}
		{$y<1$}
		{$y>2$}
		{$y<2$}}
	{\begin{tikzpicture}[scale=.7,thick,>=stealth,declare function={xmin=-2;xmax=4;ymin=-1.5;ymax=3.5;m=1;}]
			\fill[pattern=north east lines] (xmin,ymin) rectangle (xmax,ymax);
			\fill[white] (xmin,m) rectangle (xmax,ymax);
			\draw[thick,dashed] (xmin,m) -- (xmax,m);
			\draw[->] (xmin,0) -- (xmax,0)node[above]{$x$};
			\foreach \x in {-1,1,2,3,4}
			\draw[shift={(\x,0)}] (0pt,2pt) -- (0pt,-2pt) node[below] {\footnotesize $\x$};
			\draw[->,color=black] (0,ymin) -- (0,ymax)node[left]{$y$};
			\foreach \y in {-1,1,2,3}
			\draw[shift={(0,\y)}] (2pt,0pt) -- (-2pt,0pt) node[left] {\footnotesize $\y$};
			\node[above left] at (0,0){$O$};
		\end{tikzpicture}}
	\loigiai{
		Đường thẳng bờ $d\colon y=1$ vẽ nét đứt nên không kể bờ.\\
		Lấy điểm $M(0;2)$ thuộc phần không bị gạch, ta có $2>1$ nên phần không bị gạch là miền nghiệm của bất phương trình $y>1$.
	}
\end{ex}

%%%==============Câu_8 — y = 2, tô phía trên, bờ NÉT LIỀN=============%%%
\begin{ex}%[VNPmath - tự biên soạn]%[GV. Vũ Ngọc Phát]%[0D2N1-2]
	\immini[thm]{Phần không bị gạch trong hình vẽ là miền nghiệm của bất phương trình nào sau đây?
		\choice
		{$y\geq 2$}
		{$y\geq -2$}
		{\True $y\leq 2$}
		{$y\leq -2$}}
	{\begin{tikzpicture}[scale=.7,thick,>=stealth,declare function={xmin=-2;xmax=4;ymin=-0.5;ymax=4.5;m=2;}]
			\fill[pattern=north east lines] (xmin,ymin) rectangle (xmax,ymax);
			\fill[white] (xmin,ymin) rectangle (xmax,m);
			\draw[thick] (xmin,m) -- (xmax,m);
			\draw[->] (xmin,0) -- (xmax,0)node[above]{$x$};
			\foreach \x in {-1,1,2,3,4}
			\draw[shift={(\x,0)}] (0pt,2pt) -- (0pt,-2pt) node[below] {\footnotesize $\x$};
			\draw[->,color=black] (0,ymin) -- (0,ymax)node[left]{$y$};
			\foreach \y in {1,2,3,4}
			\draw[shift={(0,\y)}] (2pt,0pt) -- (-2pt,0pt) node[left] {\footnotesize $\y$};
			\node[above left] at (0,0){$O$};
		\end{tikzpicture}}
	\loigiai{
		Đường thẳng bờ $d\colon y=2$ (vẽ nét liền nên kể cả bờ).\\
		Lấy điểm $O(0;0)$ thuộc phần không bị gạch, ta có $0\leq 2$ nên phần không bị gạch là miền nghiệm của bất phương trình $y\leq 2$.
	}
\end{ex}

%%%==============Câu_9 — y = -1, tô phía trên, bờ NÉT ĐỨT=============%%%
\begin{ex}%[VNPmath - tự biên soạn]%[GV. Vũ Ngọc Phát]%[0D2N1-2]
	\immini[thm]{Phần không bị gạch (không kể đường thẳng $d$) trong hình vẽ là miền nghiệm của bất phương trình nào sau đây?
		\choice
		{$y>-1$}
		{\True $y<-1$}
		{$y>-2$}
		{$y<-2$}}
	{\begin{tikzpicture}[scale=.7,thick,>=stealth,declare function={xmin=-2;xmax=4;ymin=-3.5;ymax=1.5;m=-1;}]
			\fill[pattern=north east lines] (xmin,ymin) rectangle (xmax,ymax);
			\fill[white] (xmin,ymin) rectangle (xmax,m);
			\draw[thick,dashed] (xmin,m) -- (xmax,m);
			\draw[->] (xmin,0) -- (xmax,0)node[above]{$x$};
			\foreach \x in {-1,1,2,3,4}
			\draw[shift={(\x,0)}] (0pt,2pt) -- (0pt,-2pt) node[below] {\footnotesize $\x$};
			\draw[->,color=black] (0,ymin) -- (0,ymax)node[left]{$y$};
			\foreach \y in {-3,-2,-1,1}
			\draw[shift={(0,\y)}] (2pt,0pt) -- (-2pt,0pt) node[left] {\footnotesize $\y$};
			\node[above left] at (0,0){$O$};
		\end{tikzpicture}}
	\loigiai{
		Đường thẳng bờ $d\colon y=-1$ vẽ nét đứt nên không kể bờ.\\
		Lấy điểm $M(0;-2)$ thuộc phần không bị gạch, ta có $-2<-1$ nên phần không bị gạch là miền nghiệm của bất phương trình $y<-1$.
	}
\end{ex}

%%%%%%===================================================================%%%%%%
%%%%%%============================= DẠNG 3 ===============================%%%%%%
%%%%%% Bờ là đường thẳng ax+by = c với ab < 0  (hệ số góc dương)
%%%%%%===================================================================%%%%%%
% \subsubsection{Bờ là đường thẳng $ax+by=c$ với $ab<0$}

%%%==============Câu_10 — y = x (x - y = 0), tô nửa trên - trái, bờ NÉT LIỀN=============%%%
\begin{ex}%[VNPmath - tự biên soạn]%[GV. Vũ Ngọc Phát]%[0D2H1-2]
	\immini[thm]{Phần không bị gạch trong hình vẽ là miền nghiệm của bất phương trình nào sau đây?
		\choice
		{$y\geq x$}
		{$y\geq -x$}
		{$y\leq -x$}
		{\True $y\leq x$}}
	{\begin{tikzpicture}[scale=.7,thick,>=stealth,declare function={xmin=-2.5;xmax=3.5;ymin=-2;ymax=3;f(\x)=\x;}]
			\begin{scope}
				\clip (xmin,ymin) rectangle (xmax,ymax);
				\fill[pattern=north east lines] (xmin,ymin) rectangle (xmax,ymax);
				\fill[white] (xmin,{f(xmin)}) -- (xmax,{f(xmax)}) -- (xmax,ymin) -- (xmin,ymin) -- cycle;
				\draw[thick] (xmin,{f(xmin)}) -- (xmax,{f(xmax)});
			\end{scope}
			\draw[->] (xmin,0) -- (xmax,0)node[above]{$x$};
			\foreach \x in {-2,-1,1,2,3}
			\draw[shift={(\x,0)}] (0pt,2pt) -- (0pt,-2pt) node[below] {\footnotesize $\x$};
			\draw[->,color=black] (0,ymin) -- (0,ymax)node[left]{$y$};
			\foreach \y in {-1,1,2}
			\draw[shift={(0,\y)}] (2pt,0pt) -- (-2pt,0pt) node[left] {\footnotesize $\y$};
			\node[above left] at (0,0){$O$};
		\end{tikzpicture}}
	\loigiai{
		Đường thẳng bờ đi qua gốc toạ độ $O$ và điểm $(1;1)$ nên có phương trình $y=x$ (vẽ nét liền nên kể cả bờ).\\
		Lấy điểm $M(1;0)$ thuộc phần không bị gạch, ta có $0\leq 1$ nên phần không bị gạch là miền nghiệm của bất phương trình $y\leq x$.
	}
\end{ex}

%%%==============Câu_11 — y = 2x (2x - y = 0), tô nửa trên - trái, bờ NÉT LIỀN=============%%%
\begin{ex}%[VNPmath - tự biên soạn]%[GV. Vũ Ngọc Phát]%[0D2H1-2]
	\immini[thm]{Phần không bị gạch trong hình vẽ là miền nghiệm của bất phương trình nào sau đây?
		\choice
		{$y\geq 2x$}
		{$2y\geq x$}
		{$2y\leq x$}
		{\True $y\leq 2x$}}
	{\begin{tikzpicture}[scale=.7,thick,>=stealth,declare function={xmin=-2.5;xmax=3.5;ymin=-1.5;ymax=3.5;f(\x)=2*(\x);}]
			\begin{scope}
				\clip (xmin,ymin) rectangle (xmax,ymax);
				\fill[pattern=north east lines] (xmin,ymin) rectangle (xmax,ymax);
				\fill[white] (xmin,{f(xmin)}) -- (xmax,{f(xmax)}) -- (xmax,ymin) -- (xmin,ymin) -- cycle;
				\draw[thick] (xmin,{f(xmin)}) -- (xmax,{f(xmax)});
			\end{scope}
			\draw[->] (xmin,0) -- (xmax,0)node[above]{$x$};
			\foreach \x in {-2,-1,1,2,3}
			\draw[shift={(\x,0)}] (0pt,2pt) -- (0pt,-2pt) node[below] {\footnotesize $\x$};
			\draw[->,color=black] (0,ymin) -- (0,ymax)node[left]{$y$};
			\foreach \y in {-1,1,2,3}
			\draw[shift={(0,\y)}] (2pt,0pt) -- (-2pt,0pt) node[left] {\footnotesize $\y$};
			\node[above left] at (0,0){$O$};
		\end{tikzpicture}}
	\loigiai{
		Đường thẳng bờ đi qua gốc toạ độ $O$ và điểm $(1;2)$ nên có phương trình $y=2x$ (vẽ nét liền nên kể cả bờ).\\
		Lấy điểm $M(1;0)$ thuộc phần không bị gạch, ta có $0\leq 2$ nên phần không bị gạch là miền nghiệm của bất phương trình $y\leq 2x$.
	}
\end{ex}

%%%==============Câu_12 — y = x/2 + 1 (x - 2y + 2 = 0), tô nửa dưới - phải, bờ NÉT ĐỨT=============%%%
\begin{ex}%[VNPmath - tự biên soạn]%[GV. Vũ Ngọc Phát]%[0D2H1-2]
	\immini[thm]{Phần không bị gạch (không kể đường thẳng $d$) trong hình vẽ là miền nghiệm của bất phương trình nào sau đây?
		\choice
		{\True $x-2y+2<0$}
		{$x-2y+2>0$}
		{$2x-y+2<0$}
		{$2x-y+2>0$}}
	{\begin{tikzpicture}[scale=.7,thick,>=stealth,declare function={xmin=-3;xmax=3;ymin=-1.5;ymax=3.5;f(\x)=1+(\x)/2;}]
			\begin{scope}
				\clip (xmin,ymin) rectangle (xmax,ymax);
				\fill[pattern=north east lines] (xmin,ymin) rectangle (xmax,ymax);
				\fill[white] (xmin,{f(xmin)}) -- (xmax,{f(xmax)}) -- (xmax,ymax) -- (xmin,ymax) -- cycle;
				\draw[thick,dashed] (xmin,{f(xmin)}) -- (xmax,{f(xmax)});
			\end{scope}
			\draw[->] (xmin,0) -- (xmax,0)node[above]{$x$};
			\foreach \x in {-2,-1,1,2,3}
			\draw[shift={(\x,0)}] (0pt,2pt) -- (0pt,-2pt) node[below] {\footnotesize $\x$};
			\draw[->,color=black] (0,ymin) -- (0,ymax)node[left]{$y$};
			\foreach \y in {-1,1,2,3}
			\draw[shift={(0,\y)}] (2pt,0pt) -- (-2pt,0pt) node[left] {\footnotesize $\y$};
			\node[above left] at (0,0){$O$};
		\end{tikzpicture}}
	\loigiai{
		Đường thẳng bờ $d$ đi qua hai điểm $(0;1)$ và $(2;2)$ nên có phương trình $y=\dfrac{x}{2}+1$, hay $x-2y+2=0$.\\
		$d$ vẽ nét đứt nên không kể bờ.\\
		Lấy điểm $M(0;2)$ thuộc phần không bị gạch, ta có $0-2\cdot 2+2=-2<0$ nên phần không bị gạch là miền nghiệm của bất phương trình $x-2y+2<0$.
	}
\end{ex}

%%%==============Câu_13 — y = x + 2 (x - y + 2 = 0), tô nửa trên - trái, bờ NÉT ĐỨT=============%%%
\begin{ex}%[VNPmath - tự biên soạn]%[GV. Vũ Ngọc Phát]%[0D2H1-2]
	\immini[thm]{Phần không bị gạch (không kể đường thẳng $d$) trong hình vẽ là miền nghiệm của bất phương trình nào sau đây?
		\choice
		{$x-y+2<0$}
		{\True $x-y+2>0$}
		{$x+y+2<0$}
		{$x+y+2>0$}}
	{\begin{tikzpicture}[scale=.7,thick,>=stealth,declare function={xmin=-4;xmax=2;ymin=-1.5;ymax=3.5;f(\x)=(\x)+2;}]
			\begin{scope}
				\clip (xmin,ymin) rectangle (xmax,ymax);
				\fill[pattern=north east lines] (xmin,ymin) rectangle (xmax,ymax);
				\fill[white] (xmin,{f(xmin)}) -- (xmax,{f(xmax)}) -- (xmax,ymin) -- (xmin,ymin) -- cycle;
				\draw[thick,dashed] (xmin,{f(xmin)}) -- (xmax,{f(xmax)});
			\end{scope}
			\draw[->] (xmin,0) -- (xmax,0)node[above]{$x$};
			\foreach \x in {-3,-2,-1,1,2}
			\draw[shift={(\x,0)}] (0pt,2pt) -- (0pt,-2pt) node[below] {\footnotesize $\x$};
			\draw[->,color=black] (0,ymin) -- (0,ymax)node[left]{$y$};
			\foreach \y in {-1,1,2,3}
			\draw[shift={(0,\y)}] (2pt,0pt) -- (-2pt,0pt) node[left] {\footnotesize $\y$};
			\node[above left] at (0,0){$O$};
		\end{tikzpicture}}
	\loigiai{
		Đường thẳng bờ $d$ đi qua hai điểm $(-2;0)$ và $(0;2)$ nên có phương trình $y=x+2$, hay $x-y+2=0$.\\
		$d$ vẽ nét đứt nên không kể bờ.\\
		Lấy điểm $M(0;1)$ thuộc phần không bị gạch, ta có $0-1+2=1>0$ nên phần không bị gạch là miền nghiệm của bất phương trình $x-y+2>0$.
	}
\end{ex}

%%%%%%===================================================================%%%%%%
%%%%%%============================= DẠNG 4 ===============================%%%%%%
%%%%%% Bờ là đường thẳng ax+by = c với ab > 0  (hệ số góc âm)
%%%%%%===================================================================%%%%%%
% \subsubsection{Bờ là đường thẳng $ax+by=c$ với $ab>0$}

%%%==============Câu_14 — x + y = 1, tô nửa dưới - trái, bờ NÉT ĐỨT=============%%%
\begin{ex}%[VNPmath - tự biên soạn]%[GV. Vũ Ngọc Phát]%[0D2H1-2]
	\immini[thm]{Phần không bị gạch (không kể đường thẳng $d$) trong hình vẽ là miền nghiệm của bất phương trình nào sau đây?
		\choice
		{$x+y<1$}
		{$x-y>1$}
		{\True $x+y>1$}
		{$x-y<1$}}
	{\begin{tikzpicture}[scale=.7,thick,>=stealth,declare function={xmin=-2.5;xmax=3.5;ymin=-2;ymax=3;f(\x)=1-(\x);}]
			\begin{scope}
				\clip (xmin,ymin) rectangle (xmax,ymax);
				\fill[pattern=north east lines] (xmin,ymin) rectangle (xmax,ymax);
				\fill[white] (xmin,{f(xmin)}) -- (xmax,{f(xmax)}) -- (xmax,ymax) -- (xmin,ymax) -- cycle;
				\draw[thick,dashed] (xmin,{f(xmin)}) -- (xmax,{f(xmax)});
			\end{scope}
			\draw[->] (xmin,0) -- (xmax,0)node[above]{$x$};
			\foreach \x in {-2,-1,1,2,3}
			\draw[shift={(\x,0)}] (0pt,2pt) -- (0pt,-2pt) node[below] {\footnotesize $\x$};
			\draw[->,color=black] (0,ymin) -- (0,ymax)node[left]{$y$};
			\foreach \y in {-1,1,2}
			\draw[shift={(0,\y)}] (2pt,0pt) -- (-2pt,0pt) node[left] {\footnotesize $\y$};
			\node[above left] at (0,0){$O$};
		\end{tikzpicture}}
	\loigiai{
		Đường thẳng bờ $d$ đi qua hai điểm $(1;0)$ và $(0;1)$ nên có phương trình $x+y=1$.\\
		$d$ vẽ nét đứt nên không kể bờ.\\
		Lấy điểm $M(2;0)$ thuộc phần không bị gạch, ta có $2+0=2>1$ nên phần không bị gạch là miền nghiệm của bất phương trình $x+y>1$.
	}
\end{ex}

%%%==============Câu_15 — x + y = 3, tô nửa dưới - trái, bờ NÉT LIỀN=============%%%
\begin{ex}%[VNPmath - tự biên soạn]%[GV. Vũ Ngọc Phát]%[0D2H1-2]
	\immini[thm]{Phần không bị gạch trong hình vẽ là miền nghiệm của bất phương trình nào sau đây?
		\choice
		{$x+y<3$}
		{\True $x+y\geq 3$}
		{$x-y\geq 3$}
		{$x+y\leq 3$}}
	{\begin{tikzpicture}[scale=.7,thick,>=stealth,declare function={xmin=-1.5;xmax=4.5;ymin=-1;ymax=4;f(\x)=3-(\x);}]
			\begin{scope}
				\clip (xmin,ymin) rectangle (xmax,ymax);
				\fill[pattern=north east lines] (xmin,ymin) rectangle (xmax,ymax);
				\fill[white] (xmin,{f(xmin)}) -- (xmax,{f(xmax)}) -- (xmax,ymax) -- (xmin,ymax) -- cycle;
				\draw[thick] (xmin,{f(xmin)}) -- (xmax,{f(xmax)});
			\end{scope}
			\draw[->] (xmin,0) -- (xmax,0)node[above]{$x$};
			\foreach \x in {-1,1,2,3,4}
			\draw[shift={(\x,0)}] (0pt,2pt) -- (0pt,-2pt) node[below] {\footnotesize $\x$};
			\draw[->,color=black] (0,ymin) -- (0,ymax)node[left]{$y$};
			\foreach \y in {-1,1,2,3}
			\draw[shift={(0,\y)}] (2pt,0pt) -- (-2pt,0pt) node[left] {\footnotesize $\y$};
			\node[above left] at (0,0){$O$};
		\end{tikzpicture}}
	\loigiai{
		Đường thẳng bờ đi qua hai điểm $(3;0)$ và $(0;3)$ nên có phương trình $x+y=3$ (vẽ nét liền nên kể cả bờ).\\
		Lấy điểm $M(4;0)$ thuộc phần không bị gạch, ta có $4+0=4\geq 3$ nên phần không bị gạch là miền nghiệm của bất phương trình $x+y\geq 3$.
	}
\end{ex}

%%%==============Câu_16 — x + 2y = 2, tô nửa trên - phải, bờ NÉT ĐỨT=============%%%
\begin{ex}%[VNPmath - tự biên soạn]%[GV. Vũ Ngọc Phát]%[0D2H1-2]
	\immini[thm]{Phần không bị gạch (không kể đường thẳng $d$) trong hình vẽ là miền nghiệm của bất phương trình nào sau đây?
		\choice
		{$x+2y>2$}
		{$2x+y>2$}
		{\True $x+2y<2$}
		{$2x+y<2$}}
	{\begin{tikzpicture}[scale=.7,thick,>=stealth,declare function={xmin=-2;xmax=4;ymin=-2;ymax=3;f(\x)=1-(\x)/2;}]
			\begin{scope}
				\clip (xmin,ymin) rectangle (xmax,ymax);
				\fill[pattern=north east lines] (xmin,ymin) rectangle (xmax,ymax);
				\fill[white] (xmin,{f(xmin)}) -- (xmax,{f(xmax)}) -- (xmax,ymin) -- (xmin,ymin) -- cycle;
				\draw[thick,dashed] (xmin,{f(xmin)}) -- (xmax,{f(xmax)});
			\end{scope}
			\draw[->] (xmin,0) -- (xmax,0)node[above]{$x$};
			\foreach \x in {-1,1,2,3,4}
			\draw[shift={(\x,0)}] (0pt,2pt) -- (0pt,-2pt) node[below] {\footnotesize $\x$};
			\draw[->,color=black] (0,ymin) -- (0,ymax)node[left]{$y$};
			\foreach \y in {-1,1,2}
			\draw[shift={(0,\y)}] (2pt,0pt) -- (-2pt,0pt) node[left] {\footnotesize $\y$};
			\node[above left] at (0,0){$O$};
		\end{tikzpicture}}
	\loigiai{
		Đường thẳng bờ $d$ đi qua hai điểm $(0;1)$ và $(2;0)$ nên có phương trình $y=1-\dfrac{x}{2}$, hay $x+2y=2$.\\
		$d$ vẽ nét đứt nên không kể bờ.\\
		Lấy điểm $O(0;0)$ thuộc phần không bị gạch, ta có $0+2\cdot 0=0<2$ nên phần không bị gạch là miền nghiệm của bất phương trình $x+2y<2$.
	}
\end{ex}
\begin{ex}%[Ninh Tiến Nam - Dự án 2019-THPTMaT-1]%[0D4K4-2]%
	Cho hệ bất phương trình $\begin{cases}
			x+2y-3<0 \\
			2x+y-2>0
		\end{cases}$. Điểm nào sau đây thuộc miền nghiệm của hệ bất phương trình đã cho?
	\choice
	{$Q(-1;-5)$}
	{$N(2;2) $}
	{\True$P(3;-1) $}
	{$M(2;3)$}
	\loigiai{
		Thay tọa độ từng điểm vào ta thấy chỉ có điểm $P(3;-1)$ thỏa hai bất phương trình của hệ.}
\end{ex}

\begin{ex}%[0D4K4-4]
	\immini[thm]{Hình vẽ dưới đây là biểu diễn hình học tập nghiệm của hệ bất phương trình nào? (với miền nghiệm là miền \textbf{không} gạch sọc và chứa bờ)
		\choice
		{\True $\heva{&3x+4y-8 \geq 0 \\ &5x-12y-3\leq 0.}$}
		{$\heva{&3x+4y-8 \leq 0 \\ &5x-12y-3\leq 0.}$}
		{$\heva{&3x+4y-8 \geq 0 \\ &5x-12y-3 \geq 0.}$}
		{$\heva{&3x+4y-3 \geq 0 \\ &5x-12y-8 \leq 0.}$}}{
		\begin{tikzpicture}[line join=round, line cap=round, >=stealth,font=\footnotesize, scale=0.6]
			\draw[->](-3,0)--(4,0) node[below right] {$x$};
			\draw[->](0,-2)--(0,4) node[above] {$y$};
			\node (0,0) [above left]{$ O $};
			\foreach \x in {-2,...,2}
			\draw[shift={(\x,0)},color=black] (0pt,2pt) -- (0pt,-2pt);
			\foreach \y in {-1,...,3}
			\draw[shift={(0,\y)},color=black] (2pt,0pt) -- (-2pt,0pt);
			\draw[samples=100,smooth,domain=-2:4] plot(\x,{-0.75*(\x)+2});
			\draw[samples=100,smooth,domain=-3:4] plot(\x,{5/12*(\x)-1/4});
			\draw[dashed] (3,0)node[below] {$3$}--(3,1)--(0,1)node[left] {$1$};
			\fill[pattern=north east lines,pattern color=orange] (-2,3.5)--(-3,3.5)--(-3,-2)--(4,-2)--(4,-1);
			\fill[pattern=dots,pattern color=blue] (-3,-1.5)--(-3,-2)--(4,-2)--(4,1.416666);
			\node[right] at (0,2.1) {$2$};
		\end{tikzpicture}}
	\loigiai{Xét các điểm $A(0; 2)$, $B (3; 1)$ và $C \left(0; - \dfrac{1}{4}\right)$ thuộc bờ. \\
		Điểm $B (3; 1)$ không thỏa mãn bất phương trình $3x+4y-8 \leq 0$ nên loại $\heva{&3x+4y-8 \leq 0 \\ &5x-12y-3\leq 0.}$ \\
		Điểm $A (0; 2)$ không thỏa mãn bất phương trình $5x - 12 y - 3 \geq 0$ nên loại $\heva{&3x+4y-8 \geq 0 \\ &5x-12y-3 \geq 0.}$\\
		Điểm $C \left(0; - \dfrac{1}{4}\right)$ không thỏa mãn bất phương trình $3x + 4y - 3 \geq 0$ nên loại $\heva{&3x+4y-3 \geq 0 \\ &5x-12y-8 \leq 0.}$
	}
\end{ex}

\begin{ex}%[0D4K4-2]
	Xét $x$, $y$ thỏa mãn hệ bất phương trình $\heva{&x\ge 0 \\ & y\ge 0 \\& x+2y\le 4 \\& x-y\le 1.}$ Tìm giá trị lớn nhất $M$ của biểu thức $F=3x+2y$.\\
	% \choice
	% {\True $M=8$}
	% {$M=10$}
	% {$M=6$}
	% {$M=9$}
	\textbf{Lời giải.}
	\immini{Miền nghiệm của hệ bất phương trình là miền tứ giác $OABC$ với $O(0;0)$, $A(......)$, $B(......)$, $C(......)$.\\
	Thay toạ độ các điểm $O$, $A$, $B$, $C$ vào biểu thức ${F=3x+2y}$ ta được các giá trị tương ứng.\\
	$F(0;0)=0$, $F(......)=...$, $F(......)=...$, $F(......)=...$.\\
	Vậy giá trị lớn nhất của $F(x;y)$ bằng $M=...$ khi ${(x;y)=(...;...)}$.}
	{\begin{tikzpicture}[scale=0.8, font=\footnotesize, line join=round, line cap=round, >=stealth]
				\def\xmin{-1} \def\xmax{5}
				\def\ymin{-1} \def\ymax{3}
				\clip(\xmin,\ymin) rectangle (\xmax,\ymax);
				\tkzDefPoints{0/0/O,1/0/M,2/1/N,0/2/P}
				% \tkzDrawPoints(O,M,N,P)
				% \fill[pattern=north east lines,pattern color=blue!60] (O)--(M)--(N)--(P)--cycle;
				% \draw[domain=-1:5] plot(\x,{2-(\x)/2}) plot(\x,{(\x)-1});
				\begin{scriptsize}
					\draw[->](\xmin,0)--(\xmax,0); \draw(\xmax-0.1,0) node[below]{$x$};
					\draw[->](0,\ymin)--(0,\ymax); \draw(0,\ymax-0.2) node[right]{$y$};
					\foreach \x in {1,...,4}
					\draw (\x,0.05) -- ++(0,-0.1) node [below] {$\x$};
					\foreach \x in {1,2}
					\draw (0.05,\x) -- ++(-0.1,0) node [left] {$\x$};
					\draw node [below left]{$O$};
				\end{scriptsize}
			\end{tikzpicture}}
	\loigiai{
		\immini{
			Miền nghiệm là tứ giác như hình vẽ.\\
			$F$ lớn nhất là $8$ tại đỉnh $(2;1)$.
		}
		{
			\begin{tikzpicture}[scale=0.7, font=\footnotesize, line join=round, line cap=round, >=stealth]
				\def\xmin{-1} \def\xmax{5}
				\def\ymin{-1.5} \def\ymax{3}
				\clip(\xmin,\ymin) rectangle (\xmax,\ymax);
				\tkzDefPoints{0/0/O,1/0/M,2/1/N,0/2/P}
				\tkzDrawPoints(O,M,N,P)
				\fill[pattern=north east lines,pattern color=blue!60] (O)--(M)--(N)--(P)--cycle;
				\draw[domain=-1:5] plot(\x,{2-(\x)/2}) plot(\x,{(\x)-1});
				\begin{scriptsize}
					\draw[->](\xmin,0)--(\xmax,0); \draw(\xmax-0.1,0) node[below]{$x$};
					\draw[->](0,\ymin)--(0,\ymax); \draw(0,\ymax-0.2) node[right]{$y$};
					\foreach \x in {1,2,4}
					\draw (\x,0.05) -- ++(0,-0.1) node [below] {$\x$};
					\foreach \x in {1,2}
					\draw (0.05,\x) -- ++(-0.1,0) node [left] {$\x$};
					\draw node [below left]{$O$};
				\end{scriptsize}
			\end{tikzpicture}
		}
	}
\end{ex}

\begin{ex}%[0D4T4]
	Một cửa hàng làm kệ sách và bàn làm việc. Mỗi kệ sách cần 5 giờ chế biến gỗ và 4 giờ hoàn thiện. Mỗi bàn làm việc cần 10 giờ chế biến gỗ và 3 giờ hoàn thiện. Mỗi tháng cửa hàng có 600 giờ lao động để chế biến gỗ và 240 giờ để hoàn thiện. Lợi nhuận của mỗi kệ sách là 400 nghìn đồng và mỗi bàn là 750 nghìn đồng. Có bao nhiêu sản phẩm mỗi loại cần được làm mỗi tháng để thu được lợi nhuận tối đa?\\
	% \choice
	% {$24000$}
	% {$45000$}
	% {\True $45600$}
	% {$46000$}
	\textbf{Lời giải:}
	Ta có hệ \dongcham{4}\\
	\begin{center}
		\begin{tikzpicture}[>=stealth', scale=0.05]
				\def\xmin{-10} \def\xmax{140}
				\def\ymin{-10} \def\ymax{100}
				\clip(\xmin,\ymin) rectangle (\xmax,\ymax);
				\tkzDefPoints{\xmax/\ymax/A1,\xmin/\ymax/A2,\xmin/\ymin/A3,\xmax/\ymin/A4}
				% \fill[pattern=north east lines,pattern color=blue!60] (A1)--(A2)--(A3)--(A4)--cycle;
				\tkzDefPoints{0/0/O,0/60/A,24/48/B,60/0/C}
				% \fill[color=white] (O)--(A)--(B)--(C)--cycle;
				% \tkzDrawPoints(A,B,C)
				% \draw[domain=-10:140] plot(\x,{60-(\x)/2}) plot(\x,{80-4*(\x)/3});
				\begin{scriptsize}
					\draw[->](\xmin,0)--(\xmax,0); \draw(\xmax-4,0) node[below]{$x$};
					\draw[->](0,\ymin)--(0,\ymax); \draw(0,\ymax-4) node[right]{$y$};
					\draw node [below left]{$O$};
				\end{scriptsize}
			\end{tikzpicture}
	\end{center}
			Lợi nhuận $F(x;y)=$\dotfill\\
			Miền nghiệm của hệ là miền tứ giác $OABC$ với \dotfill \\
			Thay toạ độ các điểm $O$, $A$, $B$, $C$ vào biểu thức ${F=.............}$ ta được các giá trị tương ứng.\\
	$F(0;0)=0$, $F(.........)=......$, $F(.........)=......$, $F(.........)=......$.\\
	Suy ra giá trị lớn nhất của $F(x;y)$ bằng $.........$ khi ${(x;y)=(.....;.....)}$.\\
	Vậy \dotfill
	\loigiai{
		\immini{Ta có hệ: $\heva{& 5x+10y\le 600 \\ & 4x+3y\le 240 \\ & x\ge 0 \\ & y\ge 0}$\\
			Lợi nhuận $P=400x+750y$.\\
			Miền nghiệm của hệ là miền tứ giác $OABC$ với
			$$A(0;60),\,B(24;48),\,C(60;0)$$
			Thay tọa độ các đỉnh này vào biểu thức $P$ và so sánh kết quả, ta được
			lợi nhuận tối đa $$P_{\max}=P(B)=45600.$$
		}
		{
			\begin{tikzpicture}[>=stealth', scale=0.05]
				\def\xmin{-10} \def\xmax{140}
				\def\ymin{-10} \def\ymax{100}
				\clip(\xmin,\ymin) rectangle (\xmax,\ymax);
				\tkzDefPoints{\xmax/\ymax/A1,\xmin/\ymax/A2,\xmin/\ymin/A3,\xmax/\ymin/A4}
				\fill[pattern=north east lines,pattern color=blue!60] (A1)--(A2)--(A3)--(A4)--cycle;
				\tkzDefPoints{0/0/O,0/60/A,24/48/B,60/0/C}
				\fill[color=white] (O)--(A)--(B)--(C)--cycle;
				\tkzDrawPoints(A,B,C)
				\draw[domain=-10:140] plot(\x,{60-(\x)/2}) plot(\x,{80-4*(\x)/3});
				\begin{scriptsize}
					\draw[->](\xmin,0)--(\xmax,0); \draw(\xmax-4,0) node[below]{$x$};
					\draw[->](0,\ymin)--(0,\ymax); \draw(0,\ymax-4) node[right]{$y$};
					\draw node [below left]{$O$};
				\end{scriptsize}
			\end{tikzpicture}
		}
	}
\end{ex}
\Closesolutionfile{ans}
